\documentclass[11pt]{article}
\usepackage[T2A]{fontenc}
\usepackage[utf8]{inputenc}
\usepackage[russian]{babel}
\usepackage{lmodern}
\usepackage[margin=1in]{geometry}
\usepackage{microtype}

\emergencystretch=2em
\tolerance=1500
\hbadness=10000
\binoppenalty=700
\relpenalty=500

\usepackage{amsmath,amssymb,amsthm,mathtools,bm,mathrsfs}
\usepackage{longtable,booktabs,array,multirow}
\usepackage{textcomp}
\usepackage{graphicx}
\usepackage{enumitem}
\usepackage{xcolor}
\usepackage{framed}
\usepackage{fancyvrb}
\usepackage{ragged2e}
\usepackage{hyperref}
\usepackage{xurl}
\usepackage{tikz}
\usepackage{tikz-cd}
\usetikzlibrary{arrows.meta, positioning, calc, shapes.geometric, fit, decorations.pathreplacing}
\usepackage{caption}

% Theorem environments (Russian)
\theoremstyle{plain}
\newtheorem{theorem}{Теорема}[section]
\newtheorem{lemma}[theorem]{Лемма}
\newtheorem{proposition}[theorem]{Предложение}
\newtheorem{corollary}[theorem]{Следствие}

\theoremstyle{definition}
\newtheorem{definition}[theorem]{Определение}
\newtheorem{example}[theorem]{Пример}
\newtheorem{remark}[theorem]{Замечание}
\newtheorem{convention}[theorem]{Соглашение}

% Custom commands (same as English)
\newcommand{\RS}{\mathrm{R\text{-}S}}
\newcommand{\cS}{\mathcal{S}}
\newcommand{\cF}{\mathcal{F}}
\newcommand{\cM}{\mathcal{M}}
\newcommand{\cC}{\mathcal{C}}
\newcommand{\cD}{\mathcal{D}}
\newcommand{\cL}{\mathcal{L}}
\newcommand{\cP}{\mathcal{P}}
\newcommand{\fM}{\mathfrak{M}}
\newcommand{\Meta}{\mathfrak{Meta}}
\newcommand{\LFnd}{\mathcal{L}_{\mathrm{Fnd}}}
\newcommand{\LCls}{\mathcal{L}_{\mathrm{Cls}}}
\newcommand{\LClsMax}{\mathcal{L}_{\mathrm{Cls}}^{\top}}
\newcommand{\LAbs}{\mathcal{L}_{\mathrm{Abs}}}
\newcommand{\Sint}{\cS_{\mathrm{int}}}
\newcommand{\Ob}{\mathrm{Ob}}
\newcommand{\Mor}{\mathrm{Mor}}
\newcommand{\Hom}{\mathrm{Hom}}
\newcommand{\Fun}{\mathrm{Fun}}
\newcommand{\Cl}{\mathrm{Cl}}
\newcommand{\Lan}{\mathrm{Lan}}
\newcommand{\Ran}{\mathrm{Ran}}
\newcommand{\Set}{\mathrm{Set}}
\newcommand{\Cat}{\mathrm{Cat}}
\newcommand{\Syn}{\mathrm{Syn}}
\newcommand{\Mod}{\mathrm{Mod}}
\newcommand{\Con}{\mathrm{Con}}
\newcommand{\PA}{\mathsf{PA}}
\newcommand{\ZFC}{\mathsf{ZFC}}
\newcommand{\MLTT}{\mathsf{MLTT}}
\newcommand{\HoTT}{\mathsf{HoTT}}
\newcommand{\ETT}{\mathsf{ETT}}
\newcommand{\CIC}{\mathsf{CIC}}
\newcommand{\NBG}{\mathsf{NBG}}
\newcommand{\AFA}{\mathsf{AFA}}
\newcommand{\CZF}{\mathsf{CZF}}
\newcommand{\II}{\mathbf{I}}
\newcommand{\cU}{\mathcal{U}}
\newcommand{\tpi}{\widetilde{\pi}}
\newcommand{\cE}{\mathcal{E}}
\newcommand{\fME}{\fM_{\cE}}
\newcommand{\LAbsE}{\LAbs^{\cE}}
\newcommand{\LClsE}{\LCls^{\cE}}
\newcommand{\LClsMaxE}{(\LClsMax)^{\cE}}
\newcommand{\LFndE}{\LFnd^{\cE}}
\newcommand{\SSE}{\cS_{S}^{\cE}}
\newcommand{\Eff}{\ensuremath{\mathsf{Eff}}}
\newcommand{\NCG}{\ensuremath{\mathsf{NCG}}}
\DeclareMathOperator{\colim}{colim}
\DeclareMathOperator{\image}{image}
\DeclareMathOperator{\id}{id}
\DeclareMathOperator{\ev}{ev}
\DeclareMathOperator{\Trace}{Trace}
\DeclareMathOperator{\dep}{depth}

\title{Пространство модулей формальных систем:\\
Классификация, стабилизация и no-go-теорема для абсолютных оснований}

\author{%
  \footnotesize Макс Середа \\[-3pt]
  \scriptsize\textit{Независимый исследователь} \\[-3pt]
  \scriptsize\texttt{\href{mailto:max@gst.st}{max@gst.st}}%
}

\date{\scriptsize\today}

\begin{document}
\maketitle

\begin{abstract}
Я изучаю $(\infty, n)$-категорное пространство модулей $\fM$ формальных систем (далее \textbf{MSFS}): классифицирующий $2$-стек Морита-классов эквивалентности Rich-оснований ($\LFnd$ — теории, удовлетворяющие условиям (R1)--(R5)). Я устанавливаю пять структурных результатов и одно граничное следствие.

\emph{(i) Плюрализм в $\LCls$.} $\infty$-космои (Риль--Верити), Унивалентные основания (Воеводский) и когезивные высшие топосы (Шрайбер) — попарно $2$-неэквивалентные частичные мета-фреймворки, каждый классифицирующий собственный строгий под-стек $\fM$; они лежат вне жёсткого максимального подкласса $\Meta_{\mathrm{Cls}}^{\top}$.

\emph{(ii) Условная мета-категоричность.} Любые два члена $\Meta_{\mathrm{Cls}}^{\top}$ над одной и той же Rich-метатеорией являются $(\infty, \infty)$-эквивалентными согласно Гротендик--Люри-выпрямке с совместной точностью экстенсионального и интенсионального классификационных функторов.

\emph{(iii) Срез-локальное интенсиональное уточнение.} Функтор отображённого $2$-класса $\II : \cF^{\mathrm{op}} \to \Sint$ срез-локален над $\fM$: интенсиональные различия ($\MLTT$ против $\ETT$) расслаиваются над одиночными точками $\fM$, инвариантно разделяясь внутренне-языковой вычислимостью в эффективном топосе $\mathrm{Eff}$ Хайланда.

\emph{(iv) Мета-стабилизация на теоретическом уровне.} Итерированная мета-классификация воспроизводит ту же $(\infty, \infty)$-теорию на каждом шаге (унитарность Барвика--Шоммер-Приса), хотя теоретико-множественная инстанциация поднимается по иерархии универсумов Гротендика $\kappa_1 < \kappa_2 < \cdots$.

\emph{(v) AC/OC-двойственность.} Сопряжение Ламбека--Скотта $(\Syn, \Mod)$ из (R5b), собранное равномерно по $\RS$ и над $n \in \mathbb{N} \cup \{\infty\}$, индуцирует $(\infty, \infty)$-эквивалентность $\fM \simeq \fM_\cE$ между теоретико-центричным $2$-стеком модулей $\fM$ и актом-центричным $2$-стеком модулей $\fM_\cE$ заострённых Rich-актов. Дуал AFN-T утверждает, что акт-абсолютный стратум также пуст.

\emph{Граничное следствие: No-go-теорема для абсолютных оснований (AFN-T).} По синтаксис-семантическому сопряжению, лежащему в основе (R5), $\fM$ не имеет максимальной точки: страта $\LAbs$ (формально определимое, нередуцируемое и максимально порождающее основание) пуста. AFN-T обобщает классическую no-go-серию Кантор~\cite{Cantor1891}, Рассел~\cite{Russell1903}, Гёдель~\cite{Godel1931}, Тарский~\cite{Tarski1936}, Ловер~\cite{Lawvere1969}, Эрнст~\cite{Ernst2015} равномерно по $\RS$ и всем $(\infty, n)$, $n \in \mathbb{N} \cup \{\infty\}$, и расширяется двойственно на акт-абсолютную страту $\LAbs^\cE$.

\smallskip
\noindent\textbf{Ключевые слова:} пространство модулей формальных систем, классифицирующие $2$-стеки, категорная логика, $(\infty, n)$-категории, теория высших топосов, интенсиональная теория типов, мета-классификация, эффективный топос, no-go-теоремы.

\smallskip
\noindent\textbf{MSC2020:} 03A05, 03B30, 03C45, 03F40, 18A05, 18N10, 18N50, 18N60, 18N65.

\end{abstract}

\newpage
\tableofcontents
\newpage

\section{Введение}

Эта статья изучает \emph{пространство модулей формальных систем} $\fM$ — классифицирующий $(\infty, 2)$-стек, точки которого — Морита-классы эквивалентности Rich-оснований (формальных теорий $F$, удовлетворяющих условиям (R1)--(R5) Определения~\ref{def:rs}), а $k$-морфизмы — точные интерпретации и доказуемые естественные эквивалентности. Главные результаты касаются структуры $\fM$:
\begin{itemize}
\item \textbf{Стратификация} (\S\ref{sec:hierarchy}): $\LFnd \supsetneq \LCls \supsetneq \LClsMax \supsetneq \LAbs$.
\item \textbf{Структурный плюрализм в $\LCls$} (Теорема~\ref{thm:meta-mult}).
\item \textbf{Условная категоричность $\LClsMax$} (Теорема~\ref{thm:meta-cat}).
\item \textbf{Стабилизация на теоретическом уровне с восхождением по универсумам} (Теорема~\ref{thm:meta-stab}).
\item \textbf{Срез-локальность интенсионального уточнения} (Теорема~\ref{thm:slice-locality}).
\item \textbf{Пятиосевая абсолютность} (\S\ref{sec:five-axis}).
\item \textbf{AC/OC Морита-двойственность} (Теорема~\ref{thm:ac-oc-duality}, Теорема~\ref{thm:dual-afnt}).
\item \textbf{Внешняя граница: $\LAbs = \emptyset$} --- No-go-теорема для абсолютных оснований (Теорема~\ref{thm:afnt-alpha}, Теорема~\ref{thm:afnt-beta}).
\end{itemize}

\subsection{Обозначения и соглашения}

\begin{convention}\label{conv:zfc-inacc}
Я работаю всюду в $\ZFC$, дополненной двумя недостижимыми кардиналами $\kappa_1 < \kappa_2$, дающими универсумы Гротендика $\mathbf{U}_1 \subset \mathbf{U}_2$, достаточные для применяемой $2$-категорной техники. Категорный аппарат следует Келли~\cite{Kelly1982} для обогащённых и $2$-категорий, Люри~\cite{LurieHTT} для $(\infty, 1)$-категорий, Люри~\cite{LurieHA} для высшей алгебры, и Риль--Верити~\cite{RiehlVerity2022} для $\infty$-космоев. Для теории доступных функторов следую Адамеку--Росицкому~\cite{AdamekRosicky1994}.
\end{convention}

\begin{convention}[Ключевые символы]\label{conv:notation}
\begin{itemize}
\item $\cF$ обозначает $2$-категорию \emph{Rich-оснований} ($\LFnd$): объекты — основополагающие теории в смысле, формализованном в \S\ref{sec:rs} ниже (условия (R1)--(R5)); $1$-морфизмы — точные интерпретации; $2$-морфизмы — доказуемые естественные эквивалентности между интерпретациями.
\item $\rho : \cF \to (\infty, n)\text{-}\Cat$ обозначает \emph{классификационный функтор}, сопоставляющий каждому $F \in \cF$ его категорию моделей (или, для теоретико-типовых $F$, синтаксическую $(\infty, n)$-категорию $\Syn(F)$).
\item \emph{Калибровочное преобразование} $\tau : F_1 \to F_2$ — это $1$-морфизм в $\cF$, сопровождаемый выделенной $2$-эквивалентностью $\rho(F_1) \simeq \rho(F_2)$ на модельно-категорном уровне. Два основания \emph{калибровочно эквивалентны}, если между ними существует обратимое калибровочное преобразование.
\item $\fM$ обозначает \emph{классифицирующий $2$-стек} над $\cF$: фактор $\fM := \cF / \text{калибровка}$, рассматриваемый как $(\infty, 2)$-стек Морита-классов эквивалентности Rich-оснований.
\item Для $2$-категории $\mathbf{F}$ обозначение $\mathrm{Aut}_2(\mathbf{F})$ — $2$-группа $2$-автоэквивалентностей $\mathbf{F} \xrightarrow{\simeq_2} \mathbf{F}$.
\item $\Sint$ обозначает $2$-категорию отображённых $2$-категорий (строгое определение отнесено в Приложение~\ref{app:categorical}).
\item Все $2$-категории локально малы и $\kappa_1$-доступны, если иное не оговорено. (Для специфических Rich-оснований $S$, чья сила непротиворечивости требует второго недостижимого, параметр доступности $\Syn(S)$ обобщается до $\kappa_S \in \{\kappa_1, \kappa_2\}$ согласно условию доступности (R5a) определения Rich-метатеории.)
\end{itemize}
\end{convention}

% TODO Phase E2: §2 Stratified Hierarchy + §3 Reasonable Rich-Metatheory + §4..§7 Five-Axis Absoluteness
% TODO Phase E3: §8..§13 + Appendices

\section*{Статус перевода}

\textbf{Это пилот русского перевода (Phase E1)}. Текущий файл содержит:
\begin{itemize}
\item Полный преамбул (соглашения, theorem-окружения, символы).
\item Аннотацию (полностью переведена).
\item §1 Введение с подразделом «Обозначения и соглашения» (полностью переведено).
\end{itemize}

Остающиеся секции (§2 Стратифицированная иерархия, §3 Rich-метатеория, §4--§7 Пятиосевая абсолютность, §8--§12 ОО-вопросы и приложения) — предмет последующих фаз E2 (§2--§7) и E3 (§8--§13 + приложения).

Английский оригинал: \texttt{paper-en/paper.tex} (1692 строки, 40 страниц после Pass 20).

% Bibliography placeholder
\begin{thebibliography}{99}
\bibitem{Cantor1891} Cantor G. \"Uber eine elementare Frage der Mannigfaltigkeitslehre. \emph{Jahresbericht der DMV}, 1891.
\bibitem{Russell1903} Russell B. \emph{The Principles of Mathematics}. Cambridge UP, 1903.
\bibitem{Godel1931} G\"odel K. \"Uber formal unentscheidbare S\"atze... \emph{Monatshefte f\"ur Math.}, 1931.
\bibitem{Tarski1936} Tarski A. Der Wahrheitsbegriff in den formalisierten Sprachen. \emph{Studia Phil.}, 1936.
\bibitem{Lawvere1969} Lawvere F.W. Diagonal arguments and Cartesian closed categories. \emph{Lecture Notes Math.}, 1969.
\bibitem{Ernst2015} Ernst M. The prospects of unlimited category theory. \emph{Rev. Symb. Logic}, 2015.
\bibitem{Kelly1982} Kelly G.M. \emph{Basic Concepts of Enriched Category Theory}. CUP, 1982.
\bibitem{LurieHTT} Lurie J. \emph{Higher Topos Theory}. Princeton UP, 2009.
\bibitem{LurieHA} Lurie J. \emph{Higher Algebra}. 2017.
\bibitem{RiehlVerity2022} Riehl E., Verity D. \emph{Elements of $\infty$-Category Theory}. CUP, 2022.
\bibitem{AdamekRosicky1994} Ad\'amek J., Rosick\'y J. \emph{Locally Presentable and Accessible Categories}. CUP, 1994.
\end{thebibliography}

\end{document}
